\section{同一空间上的测度比较}
\label{sec:measure-same-space-comparison}

本节固定共同底集与可测结构$(X,\mathcal A)$。
先判断两项测度的零集是否相容，再研究一项测度的质量有多少能由相对密度表示；随后定量比较整体差异：测度之差与全变差累计正负赋权差，\term{Hellinger距离}读取平方根密度差，\term{KL散度}读取有向对数比，\term{Wasserstein距离}再额外使用共同底层距离。
最后用允许的测试函数族综合这些工具，说明它们各自能控制哪些积分读数。
这些工具按所问的问题分工：关系、表示和综合回看并不是互斥的测度类别。

零集关系本身对一般非负测度有意义；\term{Radon–Nikodym密度}和\term{Lebesgue分解}采用非负$\sigma$有限版本。
为解释两项测度之差，本小节再补有限\term{有符号测度}；其后的距离与散度统一以概率为对象。
共同空间只保证可以比较同一个事件，既不自动给出密度，也不自动提供可用的几何距离。

\subsection{\term{绝对连续}、等价与\term{相互奇异}}

本小节比较可测零集与质量集中范围。
两个子项分别处理零集的单向或双向蕴含，以及质量能否集中在互补的可测范围内：绝对连续有方向，等价要求双向共享零集，相互奇异要求可测集中范围分离。
这些关系可以分别检验，不组成互斥且穷尽的分类。

\subsubsection{绝对连续与\term{等价测度}}

\paragraph{绝对连续的定义与方向}

设$\mu,\nu$是同一可测空间上的非负测度。
先问：若$\mu$把某个可测集合视为零质量，$\nu$是否也必须如此？
若对每个$A\in\mathcal A$都有
\[
 \mu(A)=0\ \Longrightarrow\ \nu(A)=0,
\]
称$\nu$相对于$\mu$绝对连续，记为$\nu\ll\mu$。
方向写在“$\nu$相对于$\mu$”这一关系中；比较的是整个可测集合的零质量，而不是逐点密度值的大小。

若已经存在非负可测函数$f$使
\[
 \nu(A)=\int_A f\,d\mu\qquad(A\in\mathcal A),
\]
则在$\mu(A)=0$的集合上积分为零，立刻得到$\nu(A)=0$，所以这种密度表示足以推出$\nu\ll\mu$。
这只是一个方向；能否从零集蕴含反推出这样的密度函数，需要后面的Radon–Nikodym定理及其假设。

绝对连续并非连续空间专有。
实轴上相对于勒贝格测度有密度的权重满足绝对连续，\term{点质量}$\delta_0$却不满足：$\lambda(\{0\})=0$而$\delta_0(\{0\})=1$。
在有限或可数离散底集上取计数测度作参考，也可以同样定义绝对连续。
还要注意只检验单点并不够。
在实轴上令$P,Q$分别为均匀分布于$[0,1]$与$[2,3]$的概率，所有单点在两者下质量都为零，但$Q([0,1])=0$而$P([0,1])=1$，所以$P$不对$Q$绝对连续。
一般可测集合可能不可数，不能用只对可数并成立的可加性，把它的质量改写成全部单点质量之和；单点检验没有覆盖定义中的全体集合量词。

\paragraph{等价测度与共同零集}

在单向绝对连续的基础上，若$\mu\ll\nu$与$\nu\ll\mu$同时成立，记为$\mu\sim\nu$，称两测度等价；这恰好表示两者具有相同的可测零集。
等价并不意味着相等：非零集合的质量可以按不同比例分配，只是“谁质量为零”保持一致。

若已有合法表示$\nu=f\mu$且$f>0$在$\mu$几乎处处成立，则可补足反向绝对连续。
设$\nu(A)=0$，即$\int_Af\,d\mu=0$。对每个正整数$m$，由
\[
 f\,\mathbf 1_A\geq\frac1m\,\mathbf 1_{A\cap\{f\geq1/m\}}
\]
与积分单调性得到
\[
 \mu\bigl(A\cap\{f\geq1/m\}\bigr)\leq m\int_Af\,d\mu=0.
\]
正值层$\{f>0\}=\bigcup_m\{f\geq1/m\}$上$A$的部分是可数个零测集的并，质量仍为零；与$f>0$几乎处处成立合并，得到$\mu(A)=0$。
所以用几乎处处为正的密度重新加权，只改变质量大小、不改变零集。
这是“已经有密度且几乎处处为正”的充分判据；密度何时存在、何时能取倒数，将在相对密度小节补齐条件。

\subsubsection{相互奇异}

\paragraph{相互奇异与可测集中范围}

与共享零集相对的另一个问题是：两项质量能否分别集中在互补范围。
若存在可测集合$E$使
\[
 \mu(X\setminus E)=0,\qquad \nu(E)=0,
\]
称$\mu$与$\nu$相互奇异，记为$\mu\perp\nu$。
前一式表示$\mu$的全部质量集中在$E$上，后一式表示$\nu$的全部质量集中在补集$X\setminus E$上。

对勒贝格测度$\lambda$与$\delta_0$，取$E=\mathbb R\setminus\{0\}$即可：$\lambda(\{0\})=0$，$\delta_0(\mathbb R\setminus\{0\})=0$。
这里$E=\mathbb R\setminus\{0\}$与$\{0\}$不相交，但$E$的闭包是整条实轴，与$\{0\}$的闭包仍有交点。
相互奇异要求的是可测集中范围分离，并不要求两个范围的闭包不相交。

一项测度还可能同时含有相对于参考绝对连续的质量和与参考相互奇异的质量，这时它整体既不绝对连续也不相互奇异。
奇异性也不要求奇异部分由单点组成：后面的Lebesgue分解处理的是整个测度的两部分，奇异部分可以是支撑在零测的不可数集合上的连续质量。

\subsection{相对密度与测度分解}

本小节回答四个问题：密度怎样取得，密度值表示什么，怎样用它改写积分，以及哪些质量必须另外分解。
采用非负$\sigma$\term{有限测度}版本的Radon–Nikodym定理：存在结论直接陈述，不展开完整存在证明；几乎处处唯一性、单点解释和换参考公式分别说明依据。
这里的“导数”是相对积分权重，不是普通坐标导数。

\subsubsection{相对密度的存在与几乎处处唯一}

\paragraph{相对密度的存在与有限性}

绝对连续$\nu\ll\mu$只描述零集关系。
更进一步，希望用一个函数恢复$\nu$在所有可测集合上的权重。

\begin{theorem}[Radon–Nikodym定理]
\label{thm:measure-radon-nikodym}
设$\mu,\nu$是同一可测空间上的非负$\sigma$有限测度，且$\nu\ll\mu$。
则存在非负可测函数$f:X\to[0,\infty]$，使
\[
 \nu(A)=\int_A f\,d\mu\qquad(A\in\mathcal A).
\]
\end{theorem}

这个$f$记为$f=\dfrac{d\nu}{d\mu}$，称为$\nu$相对于$\mu$的Radon–Nikodym密度。
$\sigma$有限是这里所采用定理的充分假设，不写成一切密度表示的必要条件；对已有连续密度或可数离散权重的情形，可直接把恢复集合权重的已知函数辨认为Radon–Nikodym密度。

目标测度$\nu$的$\sigma$有限性还能保证密度几乎处处有限。
取可测$E_k$覆盖$X$且$\nu(E_k)<\infty$，令$N=\{f=\infty\}$。
若$\mu(N\cap E_k)>0$，则
\[
 \nu(E_k)=\int_{E_k}f\,d\mu\geq\int_{N\cap E_k}f\,d\mu=+\infty,
\]
与$\nu(E_k)<\infty$矛盾，故每个$\mu(N\cap E_k)=0$。
由可数覆盖$\mu(N)=0$，可以选取$f$在$\mu$几乎处处有限的代表，这也为后面等价测度的倒数公式准备了条件。

\paragraph{密度的几乎处处唯一}

若$f,g$都恢复$\nu$在每个可测集合上的质量，则它们只能在零集上不同。
除去二者取无穷值的共同$\mu$零集，取$\mu$有限的可测覆盖$\{M_j\}$和$\nu$有限的可测覆盖$\{N_k\}$，在每个$M_j\cap N_k$内比较；再按$f,g\leq l$的有界取值层和$f\geq g+1/m$的间隔层细分。
这些限制保证比较时积分有限，并把“严格大于”转化为至少$1/m$的正间隔。

对这样的可测层$D$，两种表示给出$\int_Df\,d\mu=\nu(D)=\int_Dg\,d\mu$。
若$\mu(D)>0$，正间隔迫使左侧比右侧至少多$\mu(D)/m$，产生矛盾；所以每个这样的层都是零测集。
可数合并各质量块、取值层和间隔层，得到$\mu(\{f>g\})=0$；交换$f,g$再得$\mu(\{g>f\})=0$。
因此
\[
 \frac{d\nu}{d\mu}=f=g\quad\mu\text{几乎处处}.
\]
在$\mu$零集上改变代表值不会改变任何集合权重；导数记号并没有指定每个点的数值，下一目就用单点质量把这一区别算清。

\subsubsection{密度值与单点质量}

\paragraph{密度系数与单点质量的区别}

固定$\nu=f\mu$并先确认$\{x\}$可测，则
\[
 \nu(\{x\})=\int_{\{x\}}f\,d\mu.
\]
这个公式比较的是单点集合的累计质量与恢复质量的函数系数，两者一般不是同一个数。

在$[0,1]$上以勒贝格测度为参考，取概率密度$p(x)=1$和$q(x)=2x$。
两项测度给每个单点都赋零质量，但在区间$[0,1/2]$上的质量分别为
\[
 \int_0^{1/2}1\,dx=\frac12,\qquad
 \int_0^{1/2}2x\,dx=\frac14.
\]
密度表示的是权重如何沿集合分布，并不是“在该点的概率”。

改换参考条件后结论才不同。
若$0<\mu(\{x\})<\infty$，单点上的积分退化为$f(x)\mu(\{x\})$，此时存在密度表示才可由
\[
 f(x)=\frac{\nu(\{x\})}{\mu(\{x\})}
\]
决定点系数，这正是有限或可数离散权重中的质量比。
若$\mu(\{x\})=0$，改变$f(x)$并不改变测度，几乎处处唯一性不会指定该点的值。

还可用$\delta_0$相对于勒贝格测度检验表示是否可能。
任何普通非负勒贝格密度$f$在$\{0\}$上的积分都为零，却应恢复质量$\delta_0(\{0\})=1$，所以不存在这样的普通可测密度。
广义函数理论中的Dirac记号需要另建分布框架，不能在这里冒充Radon–Nikodym可测函数。
缺少整个测度的参考密度，也不能靠在某个点填写无穷值补救；不能表示的奇异质量更未必只由单点构成。

\subsubsection{积分的换测度公式}

\paragraph{用相对密度改写积分}

取$f=d\nu/d\mu$。
先在指标函数$g=\mathbf 1_A$上由Radon–Nikodym表示直接得到
\[
 \int_X g\,d\nu=\nu(A)=\int_A f\,d\mu=\int_X gf\,d\mu.
\]
这是固定底层可测空间、把同一个被积对象改由另一参考测度累计，没有移动输入点，也没有重新生成输出集合。

对非负简单函数把指标恒等式有限相加；再对非负可测$g$取简单函数单调逼近，两侧分别使用单调收敛，得到
\[
 \int_X g\,d\nu=\int_X gf\,d\mu,
\]
允许两边共同为$+\infty$。
可能出现的$0\cdot\infty$按非负积分的零贡献约定处理，不在某个零密度点上强行使用普通有限数乘法。
对有符号$g$则分别计算正负部分，只在不出现$+\infty-(+\infty)$的合法范围内合并。
特别地，
\[
 \int|g|\,d\nu=\int|g|f\,d\mu
\]
使绝对可积条件在两种写法中对齐；可测性、存在密度和积分有限仍是三项不同检查。

\paragraph{重加权、限制与归一化}

换参考公式也可以反向使用：用非负可测权重$w$在$\mu$上构造新测度
\[
 \nu(A)=\int_A w\,d\mu.
\]
指标上的非负可加性与单调收敛保证$\nu$可数可加。
若归一化常数
\[
 Z=\int_Xw\,d\mu
\]
满足$0<Z<\infty$，则$Q(A)=\nu(A)/Z$是\term{概率测度}，相对于$\mu$的密度为$w/Z$；$Z=0$或$Z=+\infty$时不能按这个公式归一化，权重可测不保证$Z$有限正。

取$w=\mathbf 1_E$得到\term{限制测度}
\[
 \mu_E(A)=\mu(A\cap E);
\]
当$0<\mu(E)<\infty$时，归一化后只保留$E$内的相对权重。
底集和可测事件都没有变，被改变的是赋权；它与把输入点送到另一空间的推前不同，也与只把测度限制到较粗子$\sigma$代数不同。
在$[0,1]$上的均匀概率下取$w=2\mathbf 1_{[0,1/2]}$，则$Z=1$，新概率在$[0,1/2]$上质量为$1$、在$(1/2,1]$上质量为$0$；它对原参考绝对连续，却不再等价，因为原来正质量的区域被删去了。

\paragraph{链式换参考与倒数公式}

若$\rho\ll\nu\ll\mu$且三者都在当前非负$\sigma$有限范围内，依次对集合指标应用两次换参考公式，得到
\[
 \frac{d\rho}{d\mu}
  =\frac{d\rho}{d\nu}\frac{d\nu}{d\mu}
  \qquad\mu\text{几乎处处}.
\]
在$d\nu/d\mu=0$处把乘积的代表取为零，避免把未定的$0\cdot\infty$当成通常乘法。

若进一步$\mu\sim\nu$，令$A=\{f=0\}$，由$\nu(A)=\int_Af\,d\mu=0$与反向绝对连续得$\mu(A)=0$，故$f=d\nu/d\mu>0$在$\mu$几乎处处成立；结合前面由$\nu$的$\sigma$有限性得到的$f<\infty$，存在一个共同满测度区域使$0<f<\infty$。
在该区域定义$1/f$，对可测$B$有
\[
 \int_B\frac1f\,d\nu=\int_B\frac1f\,f\,d\mu=\mu(B),
\]
故$d\mu/d\nu=1/f$在$\nu$几乎处处成立；零集上的代表另行约定，不贡献积分。
单向绝对连续不足以取得这一步。
用简单常数例子核对方向：同一有限参考$\mu$上取$\nu=2\mu$、$\rho=3\nu$，则$d\nu/d\mu=2$、$d\rho/d\nu=3$、$d\rho/d\mu=6$，反向$d\mu/d\nu=1/2$。
公式的依据始终是恢复所有可测集合上的积分，不能形式地把微分符号约去。

\subsubsection{Lebesgue分解}

\paragraph{绝对连续部分与奇异部分}

整个测度相对于固定参考未必绝对连续，因此除了表示能写成密度的部分，还要拆出不能这样表示的部分。

\begin{theorem}[Lebesgue分解]
\label{thm:measure-lebesgue-decomposition}
设$\mu,\nu$是同一可测空间上的非负$\sigma$有限测度。
则存在唯一的非负测度$\nu_{\mathrm{ac}}$与$\nu_{\mathrm s}$，使
\[
 \nu=\nu_{\mathrm{ac}}+\nu_{\mathrm s},\qquad
 \nu_{\mathrm{ac}}\ll\mu,\qquad
 \nu_{\mathrm s}\perp\mu.
\]
\end{theorem}

对绝对连续部分应用Radon–Nikodym定理，可写成$\nu_{\mathrm{ac}}=f\mu$，于是每个可测集合满足
\[
 \nu(A)=\int_Af\,d\mu+\nu_{\mathrm s}(A).
\]
两项定理的职责要分开：分解先把测度拆成两部分，Radon–Nikodym只表示其中的绝对连续部分。
若$\nu_{\mathrm s}\neq0$，不能把$f$称为整个$\nu$相对于$\mu$的密度。
这里的唯一性针对两项测度；存在性按当前范围陈述，不展开完整证明。

\paragraph{分解的示例与参考测度依赖}

在$[0,1]$上取参考$\mu=\lambda|_{[0,1]}$及
\[
 \nu=\lambda|_{[0,1]}+\delta_0.
\]
长度部分对$\mu$绝对连续；点质量部分以$\{0\}$与$\mu$相互奇异，由此可辨认
\[
 \nu_{\mathrm{ac}}=\lambda|_{[0,1]},\qquad
 \nu_{\mathrm s}=\delta_0.
\]
集合$\{0\}$总质量为$1$，不含$0$的区间$(a,b)\subset(0,1)$总质量为$b-a$；含$0$的集合要同时累计两部分。
只有两部分相加才能恢复$\nu$对全部可测集合的权重。

分解依赖参考。保持$\nu$不变而改取$\mu'=\nu$，则$\nu$相对于$\mu'$的密度恒为$1$、奇异部分为零。
同一测度在两次分解中的分量不同，说明分解不是脱离参考的固有标签。
奇异性只要求可测集中范围分离，不要求奇异部分由点质量组成，不能把这个Dirac特例的形状推广为全部奇异测度。
下面保持共同可测空间，转向两项概率的整体差异。

\subsection{有限有符号测度与全变差}

要比较同一事件上的质量差，可以先把两项非负赋权相减。
差值允许为负，需要区分原来的非负测度、两测度之差以及它的绝对变化总量；随后再把这个通用对象用于概率之间的\term{全变差距离}。

\subsubsection{有限有符号测度与\term{总变差测度}}
\label{sec:measure-signed-total-variation}

\paragraph{集合赋权的正负分量}

对有限非负测度$\mu,\nu$，令
\[
 \eta(A)=\mu(A)-\nu(A).
\]
沿相同的不交可数分解，$\eta$仍然可数可加，但单个集合上的结果可以为负；这与“函数取正负值而积分测度仍非负”是两种不同的对象变化。
本章把取有限实值且可数可加的集合赋权称为有限有符号测度，不涉及无穷正负质量相减。

\begin{theorem}[\term{Hahn分解}与\term{Jordan分解}]
\label{thm:measure-hahn-jordan}
对有限有符号测度$\eta$，存在可测集合$E_+,E_-$满足$E_+\cup E_-=X$、$E_+\cap E_-=\varnothing$，使$\eta$在$E_+$的每个可测子集上非负、在$E_-$的每个可测子集上非正。
令
\[
 \eta^+(A)=\eta(A\cap E_+),\qquad
 \eta^-(A)=-\eta(A\cap E_-),
\]
则$\eta^+,\eta^-$是相互奇异的非负有限测度，且
\[
 \eta=\eta^+-\eta^-,
\]
这两个Jordan分量作为测度唯一。
\end{theorem}

区域$E_+,E_-$本身的选择不必唯一，但分量测度唯一。
在实轴上取$\eta=\delta_0-\delta_1$，可逐点检查$\eta(\{0\})=1$、$\eta(\{1\})=-1$、$\eta(\mathbb R)=0$，正负分量分别是$\delta_0$与$\delta_1$；全空间质量为零不代表这个集合赋权处处为零。
Jordan分解按符号分，Lebesgue分解相对于另一参考测度按绝对连续与奇异分，两者回答不同问题。

\paragraph{总变差测度与有符号积分}

定义总变差测度与\term{总变差范数}：
\[
 |\eta|=\eta^++\eta^-,\qquad
 \|\eta\|_{\mathrm{var}}=|\eta|(X).
\]
对任意可测集合$A$的有限可测分割$\{A_i\}$，都有
\[
 \sum_i|\eta(A_i)|\leq\eta^+(A)+\eta^-(A)=|\eta|(A);
\]
用$A\cap E_+$与$A\cap E_-$这两个集合分割恰好取得等号，因此$|\eta|(A)$也等于所有分割的绝对赋权和上确界。
它累计正负变化，不让两者抵消。

对满足$\int|g|\,d|\eta|<\infty$的可测$g$，定义
\[
 \int g\,d\eta=\int g\,d\eta^+-\int g\,d\eta^-,
\]
由两部分分别估计得到
\[
 \left|\int g\,d\eta\right|\leq\int|g|\,d|\eta|.
\]
有界可测函数因此满足
\[
 \left|\int g\,d\eta\right|\leq\|g\|_\infty\|\eta\|_{\mathrm{var}};
\]
控制积分要靠总变差，不能只看可能正负抵消的$\eta(X)$。

若$f$可测且$\int|f|\,d\mu<\infty$，令$\eta(A)=\int_Af\,d\mu$，则$f^+\mu$与$f^-\mu$正是它的Jordan分量，故
\[
 |\eta|=|f|\mu.
\]
这把一般有符号赋权接回已建立的非负积分。
对有限有符号测度，“所有参考零集上的赋权为零”等价于其总变差在这些零集上为零，因为可把每个参考零集再按Hahn正负区域分开检查。
本章不引入一般的向量测度理论。

\subsubsection{全变差距离与测试积分}

现在限定为同一可测空间上的概率$P,Q$。
取共同有限参考$\lambda=P+Q$，由Radon–Nikodym定理写
\[
 p=\frac{dP}{d\lambda},\qquad q=\frac{dQ}{d\lambda}.
\]
$P+Q$有限，$P,Q$都对它绝对连续，密度存在；测度之差的密度为$p-q$，总质量抵消会给出概率全变差特有的二分之一系数。

\paragraph{事件质量差与密度公式}

定义
\[
 \operatorname{TV}(P,Q)=\sup_{A\in\mathcal A}|P(A)-Q(A)|.
\]
候选是同一可测空间中的全部事件，比较的是同一事件在两种权重下的差。
由$0\leq P(A),Q(A)\leq1$可得$0\leq\operatorname{TV}\leq1$。
三角不等式由实数三角不等式逐项使用再取上确界得到；距离为零意味着所有事件质量相同，即$P=Q$；对称性由绝对值取得。
所以TV确实构成概率测度之间的距离。

在共同参考下令$h=p-q$。
总质量都为$1$给出$\int h\,d\lambda=0$，所以正部积分与负部积分相等。
取$A_+=\{p\geq q\}$，任意事件上的积分差都不超过正部总量，而$A_+$恰好达到该量。
正负两部分相加得到
\[
 \operatorname{TV}(P,Q)
  =\frac12\int|p-q|\,d\lambda
  =\frac12|P-Q|(X).
\]
二分之一来自等总质量；换另一个共同有限参考，事件差上确界不变，所以数值与参考选择无关。
推广到总质量不相等的有限测度时必须重新声明归一化约定，不能直接沿用这个$1/2$。

\paragraph{有界观测量的误差控制}

对取值于$[a,b]$的可测函数$g$，利用两项概率总质量相同，先减去常数$(a+b)/2$，再由共同参考下的密度差估计，可得
\[
 \left|\int g\,dP-\int g\,dQ\right|
  \leq(b-a)\operatorname{TV}(P,Q).
\]
特别地，$0\leq g\leq1$时系数为$1$；令$a=-\|g\|_\infty$、$b=\|g\|_\infty$，恢复一般有界函数的$2\|g\|_\infty$界。
无界观测还需要额外条件。

因为所有$0\leq g\leq1$的积分差都受TV控制，而正差集合的指标$g=\mathbf 1_{A_+}$达到上确界，所以
\[
 \operatorname{TV}(P,Q)=\sup_{0\leq g\leq1}\left|\int g\,dP-\int g\,dQ\right|.
\]
用$g=2\mathbf 1_{A_+}-1$以及有界函数界，还有
\[
 \sup_{\|g\|_\infty\leq1}\left|\int g\,dP-\int g\,dQ\right|
  =2\operatorname{TV}(P,Q).
\]
两种测试类的归一化不同，这一系数在后面比较\term{测度核}作用时还会用到。
在实轴Borel空间上，$\delta_a$与自身距离为零；$a\neq b$时事件$\{a\}$把两点区分，所以$\operatorname{TV}(\delta_a,\delta_b)=1$。
全变差辨认事件质量差异，却不读取$a,b$之间的几何间隔，这为后面的运输比较留下了明确问题。

\subsection{Hellinger距离与KL散度}

仍以同空间概率及共同有限参考$\lambda=P+Q$为对象，沿用$p=dP/d\lambda$、$q=dQ/d\lambda$。
与累计绝对密度差的全变差相比，下面分别使用平方根密度差和有向对数比，建立指标后再比较它们与全变差的关系。
这样不必在每项中重复共同参考的定义，也不把“基于密度”误写成必须有勒贝格密度。

\subsubsection{Hellinger距离}

\paragraph{平方根密度与参考无关性}

约定
\[
 H^2(P,Q)=\frac12\int(\sqrt p-\sqrt q)^2\,d\lambda,
\]
令$H(P,Q)=\sqrt{H^2(P,Q)}$。
利用$\int p\,d\lambda=\int q\,d\lambda=1$展开：
\[
 H^2(P,Q)
  =\frac12\!\left(\int p+\int q-2\int\sqrt{pq}\right)
  =1-\int\sqrt{pq}\,d\lambda.
\]
由积分形式的Cauchy–Schwarz不等式，$0\leq\int\sqrt{pq}\,d\lambda\leq1$，所以$0\leq H\leq1$。

换共同参考不改变$H$。
若改用另一共同$\sigma$有限参考$\eta$，把$\lambda,\eta$同置于$\kappa=\lambda+\eta$下，用Radon–Nikodym链式关系写出两概率相对于$\kappa$的密度；平方根差提出共同因子后再换参考积分，得到同一个$H^2$。
零密度区按积分零贡献处理，不能只凭两个公式外形相似就宣称不变。

\paragraph{距离性质与端点情形}

非负性和对称性直接由公式得到。
$H=0$意味着$\sqrt p=\sqrt q$几乎处处成立，即两密度相同，恢复的全部集合权重相同。
证明三角不等式时引入第三项概率$R$，三者统一以$P+Q+R$为参考，避免在不同参考中直接相加平方根密度；再把两段平方根差相加，对平方积分中的交叉项应用Cauchy–Schwarz即可。
所需积分不等式可由简单函数的有限和情形经逼近建立，不预先调用后章的$L^2$空间理论。

端点值可用极端情形核对：$P=Q$时$H=0$；$P\perp Q$时共同参考下两者的正密度区域不相交，$\sqrt{pq}=0$，故$H=1$。
要注意$H$和$H^2$的区别：本节称为距离的是开平方后的$H$；它与全变差、KL的关系留到\ref{sec:measure-metric-relations}小节。

\subsubsection{KL散度与有限性}

\paragraph{有向定义与积分合法性}

若$P\ll Q$，定义
\[
 D_{\mathrm{KL}}(P\|Q)=\int\log\frac{dP}{dQ}\,dP,
\]
使用自然对数；若$P$不对$Q$绝对连续，约定$D_{\mathrm{KL}}(P\|Q)=+\infty$。
在共同参考下可写成
\[
 D_{\mathrm{KL}}(P\|Q)=\int p\log\frac pq\,d\lambda,
\]
其中$p=0$的项取零；$p>0$而$q=0$的项产生$+\infty$。

令$f=dP/dQ$，则积分等于$\int f\log f\,dQ$。
由不等式
\[
 t\log t\geq-\frac1e\qquad(t\geq0)
\]
与$Q(X)=1$，被积函数的负部积分有下界，不会出现负无穷；正部仍可为$+\infty$，因此定义合法。
再由$t\log t\geq t-1$与$\int f\,dQ=1$，
\[
 D_{\mathrm{KL}}(P\|Q)=\int f\log f\,dQ\geq\int(f-1)\,dQ=0.
\]
等号要求$f=1$在$Q$几乎处处成立，即$P=Q$。
积分合法与积分有限是两件事：绝对连续不保证KL有限。

\paragraph{方向差别与无穷值的例子}

KL的方向不能互换。
对两点概率$P=(1/2,1/2)$、$Q=(3/4,1/4)$，直接代入有限和：
\[
 D_{\mathrm{KL}}(P\|Q)=\frac12\log\frac{2}{3}+\frac12\log 2
  =\frac12\log\frac43,
\]
而
\[
 D_{\mathrm{KL}}(Q\|P)
  =\frac34\log\frac32+\frac14\log\frac12,
\]
两值不同；改变方向同时改变了加权概率和对数比。
KL不对称，也不满足度量的三角不等式，不能把它当作距离使用。

即使互相绝对连续，KL也可能无穷。
在正整数上取
\[
 p_n=\frac{c}{n^2},\quad c=\Bigl(\sum_{n\geq1}n^{-2}\Bigr)^{-1},
  \qquad q_n=(e-1)e^{-n}.
\]
两者都归一化且每个点权重为正，因而互相绝对连续；但
\[
 \log\frac{p_n}{q_n}=n-2\log n+O(1),
\]
充分大的$n$上$p_n\log(p_n/q_n)$至少是某个正常数乘$1/n$，其和发散，故$D_{\mathrm{KL}}(P\|Q)=+\infty$。
这个例子把“零集关系允许定义密度”与“对数比具有足够积分控制”明确分开。

\subsubsection{全变差、Hellinger与KL的关系}
\label{sec:measure-metric-relations}

三种指标建立后，固定同一对概率、同一共同参考，判断它们能相互控制到什么程度。

\paragraph{全变差与Hellinger的双向控制}

对非负数$a,b$，有$(\sqrt a-\sqrt b)^2\leq|a-b|$，因为$|a-b|=|\sqrt a-\sqrt b|(\sqrt a+\sqrt b)$且$\sqrt a+\sqrt b\geq|\sqrt a-\sqrt b|$。
积分得到
\[
 H^2(P,Q)\leq\operatorname{TV}(P,Q).
\]
反向把$|p-q|$拆成$|\sqrt p-\sqrt q|(\sqrt p+\sqrt q)$，用Cauchy–Schwarz：
\[
 \int|p-q|\,d\lambda
 \leq\left(\int(\sqrt p-\sqrt q)^2\,d\lambda\right)^{1/2}
      \left(\int(\sqrt p+\sqrt q)^2\,d\lambda\right)^{1/2}.
\]
第二个因子不超过$2$，而$\int(\sqrt p-\sqrt q)^2\,d\lambda=2H^2$，结合$\operatorname{TV}=\frac12\int|p-q|$得到
\[
 \operatorname{TV}(P,Q)\leq\sqrt2\,H(P,Q).
\]
所以$H$与TV趋零互相等价，但数值和尺度不同。

\paragraph{Pinsker界与反向控制的边界}

自然对数约定下有\term{Pinsker不等式}
\[
 \operatorname{TV}(P,Q)^2\leq\frac12D_{\mathrm{KL}}(P\|Q).
\]
证明思路是先取全变差正差集合$A_+$，把KL按$A_+$及其补集分组；由对数函数的凸性，一般空间上的这一组不小于只在“落入$A_+$/不落入$A_+$”两个事件上的二元KL。
对二元概率，固定第二参数$b$后函数
\[
 a\mapsto a\log\frac ab+(1-a)\log\frac{1-a}{1-b}
\]
的二阶导数为$1/[a(1-a)]\geq4$，且在$a=b$处值与一阶导数为零；积分二阶下界得到该二元KL至少为$2(a-b)^2$，而$|a-b|$正是这两个事件上的全变差，整理即得Pinsker界。
端点情形由连续极限或无穷约定处理，不在边界除以零。

于是KL小先控制全变差，再控制事件概率和有界积分差；Hellinger小也能作同类控制。
反向控制一般不成立：取
\[
 P_\varepsilon=(1-\varepsilon)\delta_0+\varepsilon\delta_1,\qquad Q=\delta_0,
\]
则$\operatorname{TV}(P_\varepsilon,Q)=\varepsilon$，但$P_\varepsilon$在$\{1\}$有正质量而$Q$质量为零，故$D_{\mathrm{KL}}(P_\varepsilon\|Q)=+\infty$。
TV趋零不能推出KL趋零。

\subsection{共同底层距离下的Wasserstein比较}

现在两项概率仍在同一空间，只是比较中增加一个共同距离$d$。
先用联合权重表示两边的质量配对，再用距离幂累计配对成本；不要求预先构造把每个点送到另一个点的映射。

\subsubsection{\term{耦合}与质量配对}

\paragraph{两个边缘与耦合候选}

固定同一度量空间$X$上的Borel概率$P,Q$。
在乘积空间$X\times X$上取联合概率$\gamma$，要求
\[
 \gamma(A\times X)=P(A),\qquad
 \gamma(X\times B)=Q(B)
\]
对所有Borel集合$A,B$成立；满足这两条边缘约束的全体联合概率记为$\Pi(P,Q)$。
一个耦合同时规定从哪里取质量、与哪里的质量配对。
候选集合非空：独立耦合$P\otimes Q$总是可行。
联合权重允许把同一源点的质量分给多个目标点，因此它不是点到点函数；什么时候要求配对由确定性映射实现，会在跨空间一节讨论。

\subsubsection{Wasserstein距离}

\paragraph{距离幂成本与有限矩范围}

取带有给定完备可分度量$d$的空间$X$，本章在这一Polish范围内使用标准结论。
固定$1\leq p<\infty$和基点$x_0$，定义具有有限$p$阶矩的概率集合
\[
 \mathcal P_p(X)=\left\{P:\int_X d(x,x_0)^p\,dP(x)<\infty\right\}.
\]
换基点不改变这个集合：由幂不等式，$d(x,x_0')^p$被$d(x,x_0)^p$的常数倍加一个常数控制，反之亦然。

对$P,Q\in\mathcal P_p(X)$，定义
\[
 W_p(P,Q)
  =\left(\inf_{\gamma\in\Pi(P,Q)}
     \int_{X\times X}d(x,y)^p\,d\gamma(x,y)\right)^{1/p}.
\]
共同底层距离衡量位置差异，耦合约束保证质量守恒，开$p$次方恢复距离尺度；公式中的$x,y$都属于同一空间$X$。
至少存在一个有限成本候选：对独立耦合使用
\[
 d(x,y)^p\leq2^{p-1}\bigl[d(x,x_0)^p+d(y,x_0)^p\bigr],
\]
右侧在$P\otimes Q$下的积分由两项的有限$p$阶矩控制。
候选非空、最优成本有限和最小值由某个耦合达到是三个不同判断，不能从取下确界自动得到后两项。

\paragraph{距离性质与最优计划的范围}

在上述Polish范围与有限矩条件下，$W_p$是$\mathcal P_p(X)$上的距离，并且最优耦合存在。
三角不等式的依据是把具有共同边缘的两份耦合拼接成一份三空间耦合，再用距离的三角不等式和幂估计控制距离幂积分；同一性由$W_p(P,P)=0$与成本为零推出两边缘在对角线上一致。
存在性则要用固定边缘耦合族的紧性以及成本关于\term{弱收敛}的\term{下半连续}性，这里作为标准结果使用，不把简短说明当作完整证明。
最优联合计划不一定由确定性映射给出；若缺少有限$p$阶矩，成本可以为$+\infty$；若没有共同底层距离，也不能直接套用$d(x,y)^p$。

\paragraph{点质量与有限分布的计算}

对$\delta_a,\delta_b$，唯一耦合是$\delta_{(a,b)}$，所以
\[
 W_p(\delta_a,\delta_b)=d(a,b).
\]
这直接显示质量移动的几何距离，而不只是判断两个事件权重是否不同。

对有限离散分布
\[
 P=\sum_i\alpha_i\delta_{x_i},\qquad
 Q=\sum_j\beta_j\delta_{y_j},
\]
任何耦合都由质量矩阵$\gamma_{ij}\geq0$给出，约束为行和、列和守恒：
\[
 \sum_j\gamma_{ij}=\alpha_i,\qquad
 \sum_i\gamma_{ij}=\beta_j.
\]
在这些约束下最小化
\[
 \sum_{i,j}\gamma_{ij}d(x_i,y_j)^p,
\]
再开$p$次方就是$W_p$。
产品矩阵$\alpha_i\beta_j$给出可行点；有限维可行集闭且有界，线性成本连续，所以本例的最小值能达到。
这条有限维论证与一般\term{Polish空间}的存在性结论分别成立。

\subsection{测试函数族与比较方式的对照}

前面分别建立了测度关系、密度差异和运输成本。
本节回到比较的用途：如果最终要读取某些函数的积分，哪一种比较能控制这些读数？
测试函数族给出统一的观察口径；TV和$W_1$在其中有相应表示，但它不是与这些指标互斥的另一类分布。

\subsubsection{指定观测类的最大积分差}

\paragraph{允许的观测与可辨认范围}

固定同一可测空间上的概率$P,Q$和非空实值可测函数类$\mathcal F$，要求每个$f\in\mathcal F$在两侧都绝对可积。
把每个积分看成一种允许读取的观测，定义
\[
 d_{\mathcal F}(P,Q)
  =\sup_{f\in\mathcal F}\left|\int f\,dP-\int f\,dQ\right|,
\]
称为\term{积分概率度量}。
上确界遍历允许的全部观测，不是在源点或耦合上优化。

在共同积分可用的范围内，对称性由绝对差得到，三角不等式逐个测试成立后再取上确界。
结果可以为$+\infty$；只有在所用测度范围内有限、并且函数类能区分不同测度时，它才构成通常意义的距离。
例如在实轴上取
\[
 P=\tfrac12\delta_{-1}+\tfrac12\delta_1,\qquad Q=\delta_0,
\]
测试类只有$f(x)=x$时，两侧一阶读数都为$0$而两项测度不同，小测试类只给出伪距离。

扩大函数类会增加候选观测，上确界不会减小；若允许某个非零积分差的函数任意放大倍数，上确界就成为无穷，所以函数类的尺度限制是定义的一部分。
前面的TV公式正好是两个特例：$0\leq f\leq1$的全部可测函数给出TV，$\|f\|_\infty\leq1$给出两倍TV。
比较数值时不能忽略函数类和归一化。

\subsubsection{Lipschitz观测与一阶Wasserstein}

\paragraph{由耦合成本控制积分差}

固定Polish空间和$P,Q\in\mathcal P_1(X)$，令$g$为$L$-Lipschitz实函数。
它在两侧都可积：选定基点$x_0$后，
\[
 |g(x)|\leq|g(x_0)|+L\,d(x,x_0),
\]
有限一阶矩足以容纳这种不必有界的观测。

对任意$\gamma\in\Pi(P,Q)$，用两个边缘把积分差写成联合积分：
\[
 \int g\,dP-\int g\,dQ
  =\int[g(x)-g(y)]\,d\gamma(x,y).
\]
由Lipschitz条件$|g(x)-g(y)|\leq L\,d(x,y)$，
\[
 \left|\int g\,dP-\int g\,dQ\right|
  \leq L\int d(x,y)\,d\gamma.
\]
这对所有候选耦合成立，对右侧取下确界得到
\[
 \left|\int g\,dP-\int g\,dQ\right|
  \leq L\,W_1(P,Q).
\]
这说明运输比较怎样转化为输出误差控制；此时只证明了一个方向。

\paragraph{Lipschitz测试的对偶表示}

同一范围内还有\term{Kantorovich–Rubinstein对偶}：
\[
 W_1(P,Q)
  =\sup_{\operatorname{Lip}(f)\leq1,\ f(x_0)=0}
     \left|\int f\,dP-\int f\,dQ\right|.
\]
总质量相等允许给测试函数减去任意常数，基点归零不改变积分差，又使函数增长直接由距离控制。
反向等号是标准对偶结论，不是前一段两行估计的推论，这里不伪装成简单证明。

把它放回测试类比较：事件、有界观测和Lipschitz观测保留的信息不同。
后章的核函数空间单位球可以给出另一种测试类，但核均值、特征核条件和MMD构造依赖函数空间理论，留在那里建立；其他密度比散度和KL的变分表示也留给后面的概率分布分析，本章不穷举所有比较公式。

\subsubsection{权重差异与位置差异的对照}

固定实轴Borel空间中的$\delta_a,\delta_b$且$a\neq b$，并列核对：
\[
 \operatorname{TV}(\delta_a,\delta_b)=1,\qquad
 H(\delta_a,\delta_b)=1,\qquad
 D_{\mathrm{KL}}(\delta_a\|\delta_b)=D_{\mathrm{KL}}(\delta_b\|\delta_a)=+\infty,
\]
而
\[
 W_p(\delta_a,\delta_b)=|a-b|.
\]
前几项只使用相同事件或共同参考，最后一项还使用位置间隔；同一底集不意味着所有指标回答同一个问题。
让$b$趋向$a$但始终不重合，几何差异趋于零，固定位置的权重差异却始终为$1$：若关心事件或有界观测误差，应读取TV等控制；若允许小幅位置移动，则读取几何成本，同时核对矩条件。

再取同一Polish空间上的有限加权点概率$\widehat P=\sum_iw_i\delta_{x_i}$和无原子Borel概率$P$。
有限支撑集$S=\{x_i\}$满足$\widehat P(S)=1$、$P(S)=0$，二者相互奇异，故TV恒为$1$；有限离散近似是否有用，不能只要求它在TV下逼近连续目标。
样本或粒子是否在弱测试或几何标准下收敛，还需要后文的抽样和近似条件，本例没有自动给出一致性。
以上比较都把两项测度放在同一可测空间或同一度量空间；下一节允许输入分别来自$X,Y$，先给映射、测度核、共同表示或跨空间成本，再决定实际比较哪两个对象。
